\subsection{量子力学续}

\begin{结论}{势阱束缚态：一维总有、三维要门槛}
无限深阱 $E_n=\frac{n^2\pi^2\hbar^2}{2mL^2}$。$\delta$ 势 $V=-\alpha\delta(x)$：唯一束缚态
\[
\重点{E=-\frac{m\alpha^2}{2\hbar^2}},\qquad \psi\propto\ee^{-\kappa|x|},\ \kappa=\frac{m\alpha}{\hbar^2}\ (\psi'\text{ 跃变 }-\tfrac{2m\alpha}{\hbar^2}\psi(0)).
\]
一维任意吸引势至少一个束缚态；三维方阱（深 $V_0$、半径 $a$）需 \重点{V_0a^2>\frac{\pi^2\hbar^2}{8m}}（$s$ 波径向方程等价于一维奇宇称）。
\end{结论}

\begin{结论}{台阶与方势垒的透射}
台阶（$E>V_0$）：$R=\bigl(\frac{k_1-k_2}{k_1+k_2}\bigr)^2$，即使能量高于台阶也有反射。方势垒（高 $V_0>E$，宽 $a$，$\kappa=\sqrt{2m(V_0-E)}/\hbar$）：
\[
\重点{T=\Bigl[1+\frac{V_0^2\sinh^2\kappa a}{4E(V_0-E)}\Bigr]^{-1}}\ \xrightarrow{\kappa a\gg1}\ \frac{16E(V_0-E)}{V_0^2}\ee^{-2\kappa a}.
\]
$E>V_0$ 时把 $\sinh$ 换成 $\sin$，在 $k'a=n\pi$ 处完全透射（共振）。
\end{结论}

\begin{结论}{角动量升降与 Pauli 代数}
\[
L_\pm|l,m\rangle=\hbar\sqrt{l(l+1)-m(m\pm1)}\,|l,m\pm1\rangle,\qquad \重点{(\mathbf a\cdot\boldsymbol\sigma)(\mathbf b\cdot\boldsymbol\sigma)=\mathbf a\cdot\mathbf b+\ii\,\boldsymbol\sigma\cdot(\mathbf a\times\mathbf b)}.
\]
推论 $\ee^{\ii\theta\,\hat{\mathbf n}\cdot\boldsymbol\sigma}=\cos\theta+\ii\,\hat{\mathbf n}\cdot\boldsymbol\sigma\sin\theta$。磁场中 $H=-\gamma\mathbf S\cdot\mathbf B$：$\langle\mathbf S\rangle$ 绕 $\mathbf B$ 以 $\omega=\gamma B$ 进动，与经典磁矩相同。
\end{结论}

\begin{结论}{不确定关系与基态能量估计}
Robertson：\重点{\Delta A\,\Delta B\ge\tfrac12|\langle[A,B]\rangle|}。估计法：令 $p\sim\hbar/r$ 后对 $r$ 极小化。
氢原子 $E(r)=\frac{\hbar^2}{2mr^2}-\frac{e^2}{4\pi\varepsilon_0r}$，极小在 $r=a_0$，$E=-13.6\,$eV（巧合地精确）；谐振子得 $\hbar\omega/2$。线性势 $F|x|$：$E\sim(\hbar^2F^2/m)^{1/3}$。
\end{结论}

\begin{结论}{Hund 规则：原子基态谱项}
等价电子组态：(1) \textbf{总自旋 $S$ 最大}；(2) 在此前提下 $L$ 最大；(3) 支壳层不足半满取 $J=|L-S|$，超过半满取 $J=L+S$。
例：C $2p^2$：${}^3P_0$；N $2p^3$：${}^4S_{3/2}$；O $2p^4$：${}^3P_2$；Fe $3d^6$：${}^5D_4$。物理原因：自旋对称化使空间波函数反对称，电子彼此远离，排斥能降低（交换能）。
\end{结论}

\begin{结论}{氦原子的变分估计}
试探 $\psi\propto\ee^{-Z'(r_1+r_2)/a_0}$，$\langle H\rangle=\bigl(Z'^2-2ZZ'+\frac58Z'\bigr)$ Hartree：
\[
\重点{Z'=Z-\frac5{16}=\frac{27}{16}},\qquad E\approx-\Bigl(\frac{27}{16}\Bigr)^2\ \mathrm{Ha}=-77.5\ \mathrm{eV}\quad(\text{实验 }-79.0).
\]
$\frac{5}{16}$ 为另一电子的屏蔽。一阶微扰（$Z'=Z$）给 $-74.8\,$eV；忽略排斥给 $-108.8\,$eV。
\end{结论}

\begin{结论}{简并微扰：氢 $n=2$ 线性 Stark 效应}
外场 $\mathcal E\hat z$，$V=e\mathcal Ez$ 只连接 $2s$ 与 $2p_0$（$\Delta m=0$，宇称相反）：$\langle2s|z|2p_0\rangle=-3a_0$。
\[
\重点{\Delta E=\pm3ea_0\mathcal E},\qquad 2p_{\pm1}\ \text{不移动}.
\]
线性 Stark 效应只在简并能级间存在；非简并基态只有二阶效应 $-\frac12\alpha\mathcal E^2$，氢的极化率 $\alpha=\frac92a_0^3\cdot4\pi\varepsilon_0$。
\end{结论}

\begin{结论}{突变近似：氚 $\beta$ 衰变}
Hamilton 量突变时波函数来不及变，跃迁概率为重叠的平方。氚（$Z=1$）基态电子在核变为 $\mathrm{He}^+$（$Z=2$）后留在基态：
\[
|\langle1s_{Z_2}|1s_{Z_1}\rangle|^2=\Bigl(\frac{2\sqrt{Z_1Z_2}}{Z_1+Z_2}\Bigr)^6=\重点{\frac{512}{729}\approx0.70}.
\]
适用条件：变化时间 $\ll\hbar/\Delta E$。与之相对的绝热极限：慢变时停留在对应的瞬时本征态。
\end{结论}

\begin{结论}{分波法与光学定理}
\[
\重点{\sigma=\frac{4\pi}{k^2}\sum_l(2l+1)\sin^2\delta_l},\qquad \sigma_{\mathrm{tot}}=\frac{4\pi}{k}\Im f(0)\ (\text{光学定理}).
\]
低能只有 $s$ 波：$\delta_0\approx-ka$，$\sigma\to4\pi a^2$（$a$ 为散射长度）。硬球低能 $4\pi R^2$（经典值的 4 倍），高能 $2\pi R^2$（含衍射峰）。$l$ 波在 $ka\gtrsim l$ 时才贡献。
\end{结论}

\begin{结论}{密度矩阵、Bloch 矢量与 CHSH}
$\rho=\sum p_i|\psi_i\rangle\langle\psi_i|$，纯态 $\iff\tr\rho^2=1$。量子比特 \重点{\rho=\tfrac12(I+\mathbf r\cdot\boldsymbol\sigma)}，$|\mathbf r|\le1$，$|\mathbf r|=1$ 为纯态。
子系统约化密度矩阵的 von Neumann 熵量化纠缠（Bell 态为 $\ln2$）。CHSH 组合：局域隐变量 $|S|\le2$，量子力学最大 $2\sqrt2$（Tsirelson 界）。
\end{结论}

\begin{结论}{Berry 相位与 Aharonov--Bohm 相位}
绝热地沿闭合路径改变参数，本征态获得几何相位 $\gamma=\ii\oint\langle n|\nabla_R n\rangle\cdot\dd\mathbf R$。自旋 $\frac12$ 跟随缓慢转动的磁场：
\[
\重点{\gamma=\mp\frac12\Omega}\ (\Omega\text{ 为 }\hat{\mathbf B}\text{ 扫过的立体角}).
\]
带电粒子绕通量 $\Phi$（路径上 $\mathbf B=0$）：相位 $\frac{q\Phi}{\hbar}$，干涉条纹以通量量子 $h/e$ 为周期移动（超导中为 $h/2e$）。
\end{结论}

\begin{结论}{自由传播子与波包扩散}
\[
K(x,t;x',0)=\sqrt{\frac{m}{2\pi\ii\hbar t}}\exp\Bigl[\frac{\ii m(x-x')^2}{2\hbar t}\Bigr].
\]
初始宽度 $\sigma_0$ 的 Gauss 波包：\重点{\sigma(t)=\sigma_0\sqrt{1+\Bigl(\frac{\hbar t}{2m\sigma_0^2}\Bigr)^2}}。电子 $\sigma_0=1\,$Å 约 $10^{-16}\,$s 即显著扩散；$1\,$g 物体 $\sigma_0=1\,\mu$m 需约 $10^{19}\,$s，远超宇宙年龄。
\end{结论}

\begin{结论}{Landau 能级}
均匀磁场中二维电子：
\[
\重点{E_n=\hbar\omega_c\Bigl(n+\frac12\Bigr)},\quad \omega_c=\frac{eB}{m},\quad \text{每能级单位面积简并度 }\frac{eB}{h},\quad \ell_B=\sqrt{\frac{\hbar}{eB}}.
\]
$\ell_B\approx25.7\,$nm$/\sqrt{B/\mathrm T}$。填充因子 $\nu=\frac{nh}{eB}$；整数量子 Hall 效应中 $\sigma_{xy}=\nu\frac{e^2}{h}$，$h/e^2=25812.8\,\Omega$。
\end{结论}

\begin{结论}{Feynman--Hellmann 定理}
$H(\lambda)$ 的归一化本征态：
\[
\重点{\frac{\partial E_n}{\partial\lambda}=\Bigl\langle n\Bigm|\frac{\partial H}{\partial\lambda}\Bigm|n\Bigr\rangle}.
\]
对氢原子取 $\lambda=e^2$ 得 $\langle1/r\rangle$，取 $\lambda=l$（有效径向方程）得 $\langle1/r^2\rangle$；对谐振子取 $\lambda=\omega$ 得 $\langle V\rangle=\langle T\rangle$。分子中核受力也由它给出（Hellmann--Feynman 力）。
\end{结论}

\参考{Griffiths--Schroeter, \emph{Introduction to Quantum Mechanics}；Sakurai--Napolitano, \emph{Modern Quantum Mechanics}；Shankar, \emph{Principles of Quantum Mechanics}；Bransden--Joachain, \emph{Physics of Atoms and Molecules}；Nielsen--Chuang 第 2 章（MIT 8.05 / 8.06）。}
